% 01-intro.tex — 引言

\section{引言}
\label{sec:intro}

\subsection{问题陈述: 元胞自动机终态是否为迷宫}

元胞自动机 (Cellular Automaton, CA) 是一类离散动力系统, 其状态在网格上由局部更新规则同步演化。Conway 的 ``生命游戏'' 是最早被广泛研究的 CA 规则: 在 8 邻域下, 活细胞在邻居数属于存活集 $S$ 时继续存活, 死细胞在邻居数属于出生集 $B$ 时复活 \cite{conway1970}. Wolfram 将 CA 的动力学行为划分为四类 (均匀, 周期, 混沌, 复杂), 生命游戏属于复杂类 \cite{wolfram2002}.

CA 与迷宫生成在结构上具有相似性: 两者均产生二值网格, 依赖邻域信息, 并体现 ``局部规则涌现全局结构'' 的特征。然而, CA 动力学是双向的, 可将 ``看似迷宫'' 的图样视为 CA 的吸引子。这一观察引出一个明确的判定问题: 给定一个 CA 终态 (二值网格), 如何判断其是否构成迷宫?

该问题在直觉上似乎简单 (考察连通性与走廊结构), 但实证上至少存在四类 CA 吸引子能够通过朴素的判别:
\begin{itemize}[leftmargin=*, itemsep=0pt, topsep=0pt]
  \item \textbf{实心块}. 墙比为 1, 整张图完全黑化, 拓扑退化。
  \item \textbf{空盒}. 墙比为 0, 整张图完全白化, 同样退化。
  \item \textbf{规则条纹、同心环与蜂巢}. 高度对称的几何反例, 视觉上具有结构, 但缺乏迷宫的一格宽走廊特征。
  \item \textbf{随机噪声}. 50\% 密度的随机点可通过 ``连通性 $\ge$ 阈值'' 的检查, 但分叉过多, 不具迷宫特征。
\end{itemize}

有效拒收上述反例, 需要一个同时刻画 ``走廊结构'' 与 ``非平凡拓扑'' 的度量, 并在伪迷宫上实现判别失效。

\subsection{已有工作的两个空缺}

迷宫生成算法 (DFS \cite{pech2002}, Kruskal, Prim, Growing Tree, Sidewinder, Binary Tree) 与 CA 演化均产生二值网格, 但两者之间缺少统一的 ``是否迷宫'' 判别器。Bellot 等人 2021 年提出的 $F = \nu / \delta$ 度量 \cite{bellot2021} 是当前学界常用的对照基线。然而, 该度量对 ``少死端墙 + 长直径'' 的几何反例 (条纹、环、蜂巢) 反而给出高分, 方向与迷宫判定相反。

CA 规则搜索的现有工作多集中于经典 B/S 字符串 ($2^{18}$ 候选), 采用遗传编程或进化策略搜索 ``自然涌现迷宫'' 的规则 \cite{rejbrand2017, mcclendon2001}. 然而, 此类规则的搜索空间有限, 仅能产生有限种动力学, 难以涌现复杂走廊结构。FamilyMask 编码将搜索空间扩展至 $2^{1648}$, 同时将 B/S 字符串作为单族退化情形保留。

\subsection{本文贡献}

本文给出以下三个互补的贡献:
\begin{enumerate}[leftmargin=*, itemsep=0pt, topsep=0pt]
  \item \textbf{FamilyMask 多族 CA 染色体}. 经典 B/S 规则被推广为最多 16 族并行、按优先级仲裁的染色体表示。每族独立持有邻域 mask、出生集、存活集与优先级。搜索空间为 $2^{103}$ 单族 / $2^{1648}$ 全染色体。
  \item \textbf{maze\_quality 八维几何度量}. 由 5 个拓扑子度量与 3 个多样性子度量构成, 通过两层加权几何聚合与 $\min$ 平衡, 配合墙比三角门控, 全部连续可微, 可由 ES 直接优化。在 15-pattern 基准上, 真迷宫均值 $0.711$, 伪迷宫均值 $0.000$, 间隔 $+0.711$, 误判 $0/9$. 同数据集上 Bellot 2021 $F$ 度量间隔 $-191.3$, 方向反转, 误判 $4/9$.
  \item \textbf{两个尺度的演化扫描}. 小范围扫描 ($40 \times 60$, 122 OK) 给出完整的 (mask, maxFam) 消融数据, 最高分 $0.8233$ (manhattan-2 / maxFam=8 / seed 444). 大范围扫描 ($500 \times 2000$, 5/5 OK) 验证 ``同配置, 大网格'' 的 scale-up 性能, 最高分 $0.8195$, 5-run 均值 $0.8020$. 两扫描峰值差 $0.0037$, 表明小范围网格已接近该配置的理论上界。
\end{enumerate}

\subsection{全文结构}

§2 综述 CA 与迷宫生成、FamilyMask 类编码、几何度量与 ES 搜索等相关工作。§3 给出 FamilyMask 染色体的形式化描述与执行语义。§4 提出 maze\_quality 的 8 个子度量定义与两层聚合, 并在 15-pattern 基准上与 Bellot 2021 $F$ 进行对比。§5 报告两个尺度的扫描实验设置、结果与对比分析。§6 讨论两类失败模式 (预算受限与表征受限) 及大网格 scale-up 的含义。§7 总结全文并指出后续方向。附录 A 给出完整公式, 附录 B 给出数据可用性与复现性说明, 附录 C 给出软硬件环境。
